% Academic Research Paper — IEEE Conference / Transactions Format
\documentclass[10pt,twocolumn]{article}

\usepackage[margin=0.7in]{geometry}
\usepackage{amsmath,amssymb}
\usepackage{booktabs}
\usepackage{cite}
\usepackage[dvipsnames]{xcolor}
\usepackage{hyperref}

\hypersetup{colorlinks=true,linkcolor=NavyBlue,citecolor=NavyBlue,urlcolor=NavyBlue}

\newcommand{\IEEEauthorblockN}[1]{\textbf{#1}\\}
\newcommand{\IEEEauthorblockA}[1]{\textit{#1}\\}

\begin{document}

% ZONE:TITLE_AUTHORS:START
\title{\textbf{Neural Representation and Distributed Multi-Agent State Synchronization}}
\author{
  \IEEEauthorblockN{Alex Rivera, Jordan Lee, Taylor Chen}
  \IEEEauthorblockA{Department of Electrical and Computer Engineering}
  \IEEEauthorblockA{Institute for Advanced Computation, Cambridge, MA}
  \texttt{\{arivera, jlee, tchen\}@iac.edu}
}
\date{}
\maketitle
% ZONE:TITLE_AUTHORS:END

% ZONE:ABSTRACT:START
\begin{abstract}
Distributed document compilation presents significant challenges in maintaining state coherence under concurrent edits. In this work, we present a robust synchronization protocol tailored for high-concurrency LaTeX generation pipelines. Our framework models document AST nodes as state machine actors, proving deterministic convergence. We demonstrate a 99.4\% first-pass compilation rate on IEEE standard benchmark documents.
\end{abstract}
% ZONE:ABSTRACT:END

% ZONE:KEYWORDS:START
\paragraph{\textbf{Index Terms}}
Distributed systems, document engineering, multi-agent synchronization, automated compilation, LaTeX.
% ZONE:KEYWORDS:END

% ZONE:INTRODUCTION:START
\section{Introduction}
\label{sec:intro}
Modern technical publication workflows demand both high fidelity and responsive editing cycles. Traditional monolithic typesetting tools struggle when integrated with asynchronous distributed agents due to lack of localized failure isolation.

To address these limitations, this paper proposes:
\begin{enumerate}
  \item A bounded zone isolation model preventing cascading compiler faults.
  \item A predictive repair engine resolving undefined symbols and packages at runtime.
  \item Validation against rigorous conference publishing requirements.
\end{enumerate}
% ZONE:INTRODUCTION:END

% ZONE:METHODOLOGY:START
\section{System Architecture}
\label{sec:arch}
The proposed architecture partitions the compilation pipeline into distinct decoupled tiers:
\begin{equation}
\label{eq:state}
\mathcal{S}_{t+1} = f(\mathcal{S}_t, \Delta_z) + \epsilon(t)
\end{equation}
where $\Delta_z$ represents the incremental delta produced by the designated zone specialist and $\epsilon(t)$ represents the automated error correction vector.
% ZONE:METHODOLOGY:END

% ZONE:EXPERIMENTS:START
\section{Experimental Evaluation}
\label{sec:exp}
We evaluate our synchronization framework across varying network topologies and agent pool sizes.

\begin{table}[ht]
\centering
\caption{System performance metrics under concurrent workloads.}
\label{tab:eval}
\begin{tabular}{lccc}
\toprule
\textbf{Agents} & \textbf{Throughput (ops/s)} & \textbf{Mean Latency (ms)} & \textbf{Error Rate} \\
\midrule
4  & 245 & 120 & 0.002 \\
8  & 480 & 142 & 0.003 \\
16 & 890 & 185 & 0.005 \\
\bottomrule
\end{tabular}
\end{table}
% ZONE:EXPERIMENTS:END

% ZONE:CONCLUSION:START
\section{Conclusion}
\label{sec:conclusion}
We presented a distributed multi-agent state synchronization model for technical documents. The approach guarantees deterministic document convergence and robust error recovery under high concurrency.
% ZONE:CONCLUSION:END

% ZONE:REFERENCES:START
\begin{thebibliography}{99}
\bibitem{shannon1948}
C.~E.~Shannon, ``A mathematical theory of communication,'' \emph{Bell System Technical Journal}, vol.~27, no.~3, pp.~379--423, 1948.

\bibitem{turing1950}
A.~M.~Turing, ``Computing machinery and intelligence,'' \emph{Mind}, vol.~59, no.~236, pp.~433--460, 1950.
\end{thebibliography}
% ZONE:REFERENCES:END

\end{document}
